%================================ RÉSUMÉ =====================================
\chapter*{Résumé}
\addcontentsline{toc}{chapter}{Résumé}
\markboth{Résumé}{Résumé}

Ce projet de fin d'année porte sur la conception et le développement d'un
\textbf{prototype de plateforme web} pour \textbf{SBS Data Factory}, destiné à
la gestion de campagnes publicitaires de type \emph{drive-to-store} --- des
campagnes dont l'objectif est de convertir une audience digitale en
\textbf{visites physiques en magasin}.

La plateforme, nommée \textbf{Drive-to-Store DSP}, couvre l'ensemble du
parcours métier : authentification et navigation protégée par rôle, tableau de
bord de synthèse, gestion et création de campagnes via un assistant
multi-étapes, import et validation de la base de magasins d'un annonceur,
ciblage géographique (\emph{geofencing}) sur carte interactive, génération de
créatives dynamiques (DCO) et de \emph{landing pages} personnalisées par
magasin, reporting de performance (indicateurs, graphiques, carte, tableau et
export CSV), ainsi qu'un module d'administration du compte (annonceurs,
utilisateurs, paramètres).

Sur le plan technique, l'application repose sur un \textbf{frontend React
(Vite)} et un \textbf{backend FastAPI} exposant une API REST documentée via
\textbf{Swagger/OpenAPI}. La cartographie s'appuie sur
\textbf{Leaflet / OpenStreetMap}. Le développement a été mené de façon
\textbf{itérative} avec un flux de travail \textbf{Git/GitHub} (branches et
\emph{pull requests}), une validation fonctionnelle après chaque module et une
documentation continue.

Le livrable est un \textbf{prototype fonctionnel de bout en bout} : les
données sont, pour l'essentiel, simulées (\emph{mockées}) et conservées en
mémoire --- à l'exception des brouillons de campagne, réellement persistés
côté backend ---, l'authentification est simulée et la base PostgreSQL est
\emph{préparée mais non branchée en production}. Ces choix, assumés à ce stade, ont permis de valider l'intégralité
des parcours utilisateur avant d'envisager le branchement d'une base réelle et
d'une authentification sécurisée.

\vspace{0.6cm}
\noindent\textbf{Mots-clés :} Drive-to-Store, DSP, publicité digitale,
React, FastAPI, geofencing, DCO, reporting, prototype web, OpenStreetMap.

%================================ ABSTRACT ===================================
% Chapitre étoilé : la classe report force un changement de page, l'Abstract
% commence donc en haut d'une nouvelle page. En-tête et entrée de table des
% matières propres à l'Abstract (ne montrent jamais « Résumé »).
\chapter*{Abstract}
\addcontentsline{toc}{chapter}{Abstract}
\markboth{Abstract}{Abstract}

This end-of-year project covers the design and development of a
\textbf{web platform prototype} for \textbf{SBS Data Factory}, aimed at
managing \emph{drive-to-store} advertising campaigns --- campaigns whose goal
is to turn a digital audience into \textbf{physical in-store visits}.

The platform, called \textbf{Drive-to-Store DSP}, spans the full business
journey: authentication with role-based protected navigation, an overview
dashboard, campaign management and a multi-step creation wizard, import and
validation of an advertiser's store database, geographic targeting
(\emph{geofencing}) on an interactive map, dynamic creative generation (DCO)
with simulated per-store landing pages, performance reporting (KPIs, charts,
map, table and CSV export), and an account administration module
(advertisers, users, settings).

Technically, the application relies on a \textbf{React (Vite) frontend} and a
\textbf{FastAPI backend} exposing a REST API documented through
\textbf{Swagger/OpenAPI}, with mapping powered by
\textbf{Leaflet / OpenStreetMap}. Development followed an \textbf{iterative}
approach with a \textbf{Git/GitHub} workflow (branches and pull requests),
functional validation after each module, and continuous documentation.

The deliverable is a \textbf{fully working end-to-end prototype}. Most data is
\emph{mocked}/in-memory (except campaign drafts, which are actually persisted
on the backend), authentication is simulated, and the PostgreSQL database is
\emph{prepared but not connected in production}. These deliberate choices made
it possible to validate all user journeys before wiring a real database and
production-grade authentication.

\vspace{0.5cm}
\noindent\textbf{Keywords:} Drive-to-Store, DSP, digital advertising, React,
FastAPI, geofencing, DCO, reporting, web prototype, OpenStreetMap.
